\clearpage
\section{Introducción}
\label{iii:sec:introduccion}

Una relación constitutiva organiza la respuesta de un sistema, pero su
forma funcional y sus coeficientes plantean cuestiones diferentes. Una
vez conocidos los coeficientes, la ecuación permite calcular respuestas;
determinar por qué esos coeficientes tienen conjuntamente sus valores
requiere describir la estructura que los fija. Este trabajo aborda esa
segunda cuestión mediante una construcción matemática de dos canales
orientados. Su tesis es que la estructura precedente determina la
ecuación y que la evaluación conserva una parte recuperable de esa
estructura. El vacío electromagnético constituye aquí una realización
precisa de esa tesis: permitividad, permeabilidad, impedancia y velocidad
se obtienen como cuatro lecturas relacionadas de una misma respuesta.

El resultado principal comprende tanto la construcción directa como la
inversa. En el dominio contractivo, el operador constitutivo es positivo,
sus autovalores pertenecen a un intervalo determinado y sus cuatro
coordenadas satisfacen identidades exactas. La monotonía estricta de la
función de respuesta proporciona una inversa única; con las hojas
etiquetadas, esa inversa recupera los canales y el par angular. Esta
recuperación permite componer tres descripciones: la respuesta del
vacío, la forma de la elipse de acción y el funcional de
Barbero--Immirzi. La comparación no se limita, por tanto, a valores
numéricos: identifica las operaciones que transmiten orientación e
información espectral entre esas descripciones.

\subsection{Núcleo generador y procedencia de la respuesta}

El modelo se construye desde un núcleo formal matemático denominado
Holografía Modular Triádica, en adelante HMT. Su primer objeto,
\emph{All Possible Paths} (APP), reúne caminos y evaluaciones
aritméticas sobre las $81$ posiciones del soporte
$X_9=(\mathbb Z/9\mathbb Z)^2$. Las traslaciones de generadores
$\{(1,0),(-1,0),(0,1),(0,-1)\}$ definen su grafo de Cayley aditivo,
una retícula cíclica bidimensional con realización toroidal. Las
evaluaciones suma y producto actúan sobre ese soporte; cada una conserva
su residuo y su cociente. La estructura de caminos y las evaluaciones
tienen así funciones distintas dentro de APP: los desplazamientos
especifican las rutas, mientras las hojas aritméticas determinan los
datos que éstas transportan.

El segundo objeto, TRIT, es la firma ternaria orientada
$\{-1,0,+1\}$, obtenida mediante la representación balanceada de las
clases módulo $3$. Su realización cuadrática distingue los tipos
hiperbólico, parabólico y elíptico, respectivamente. El tercer objeto
es el \emph{Topological Phase Kernel}, en adelante TPK: compone la
selección de fase, el transporte cardinal y la actualización de memoria.
La dinámica actúa sobre posiciones, hojas, cocientes, orientación,
prefijos y datos de frontera. El retorno de una fase conserva la
historia acumulada, de modo que una vuelta puede cerrar una coordenada
sin restituir el estado completo.

La prolongación nonádica transmite esa estructura mediante prefijos y
fronteras compatibles. En su desarrollo conjunto se generan las
coordenadas $\pi_{\rm HMT}$, $\varphi_{\rm HMT}$,
$e_{\rm HMT}$ y $\alpha_{\rm HMT}$. El registro dodecafásico $K$
y su evaluación periódica intervienen en la determinación de
$\alpha$ mediante el cierre dodecafásico; la representación analítica completa expresa esa misma
coordenada en la clase de colas construida. A continuación se obtienen
las secciones de acción y el par angular. El uso de esos resultados
como antecedentes conserva sus demostraciones dentro del artículo.
Esta secuencia fija la genealogía de los canales antes de evaluar la
respuesta constitutiva.

La orientación de hojas y su incidencia areal--volumétrica cumplen
papeles complementarios. La primera distingue el retorno de una ruta
de la restauración de su orientación; la segunda procede del calendario
de modos y de su matriz de incidencia. La frecuencia cronológica, la
medida del sistema simbólico y la marca regional de frontera mantienen
sus definiciones propias. Estas distinciones, desarrolladas en la
sección~\ref{iii:sec:hojas}, precisan la información que llega al
transporte angular y evitan sustituirla por una razón escalar aislada.

\clearpage
\subsection{Determinación constitutiva y recuperación de la forma}

Con $T=q_+P_++q_-P_-$ y $0<q_\pm<1$, las incidencias temporales
$90$ y $120$ determinan
\[
\mathcal V=(I-T^{90})(I-T^{120})^{-1}.
\]
Al escribir $S=T^{30}$, la ecuación
$\mathcal V(I+S+S^2+S^3)=I+S+S^2$ muestra la regla que fija la
respuesta para ese transporte. La sección~\ref{iii:sec:constitutiva}
demuestra su existencia y unicidad y establece la función escalar
$f:(0,1)\to(3/4,1)$. Sus valores $r_\pm=f(q_\pm^{30})$
producen
\[
\widehat\varepsilon=r_+^2,\qquad
\widehat\mu=r_-^2,\qquad
\widehat Z=r_-/r_+,\qquad
\widehat c=(r_+r_-)^{-1}.
\]
La positividad y las lecturas de producto y cociente determinan los dos
cuadrados una vez fijadas las relaciones constitutivas. La unicidad
afirmada pertenece a ese sistema explícito de reglas y a su dominio.

La inversión cúbica recupera $q_\pm$ y luego las coordenadas
$x=-\frac12\log(q_+q_-)$,
$y=\frac12\log(q_-/q_+)$. Con las etiquetas de hoja conservadas,
impedancia y velocidad son suficientes para esta reconstrucción.
La proposición~\ref{iii:prop:forma-LC} obtiene a continuación la
anisotropía de la elipse de acción y la impedancia inductiva--capacitiva
relativa. Las referencias de acción, frecuencia e impedancia completan
su realización dimensional. La proposición~\ref{iii:prop:identidad-velocidad}
prueba, en la misma coordenatización, $c_{\rm int}=c_{\rm vac}$ y relaciona la
longitud de trayectoria con la duración elemental. La inversión
reconstruye así antecedentes espectrales determinados, mientras ruta
y memoria permanecen en el estado enriquecido del que proceden.

Esta inversa induce asimismo una ley de composición entre respuestas:
la multiplicación de los transportes se expresa mediante la operación
transportada y admite una coordenada aditiva. La
sección~\ref{amp-higgs-seccion} compone después los mismos canales con
las firmas angulares de calibre y de Higgs. La inversa conjunta, su
dominio exacto y la reconstrucción de impedancia conservan las condiciones
de escala, refinamiento y normalización de esa realización.

\subsection{Ciclo, constantes y funcionales orientados}

La función $f$ admite una realización sobre el ciclo de doce fases.
Sus subespacios de dimensiones $3$ y $4$ producen el cociente de
determinantes; su diferencia de traza normalizada da $-1/12$.
En un plano orientado del mismo ciclo, un cuarto de giro $J$ satisface
$J^2=-I$ y su resolvente determina $s/(1+s^2)$. Dos integraciones
logarítmicas, con sus condiciones en el origen, generan el momento
$\mathcal C_2(s)$ y su valor extremo, la constante de Catalán.
El teorema~\ref{iii:thm:catalan-pell} relaciona exactamente la
evaluación en $2-\sqrt3$ con ese valor extremo y con
$\log(2+\sqrt3)$. La sección algebraica de Pell y los argumentos
angulares $q_\pm^{30}$ son evaluaciones distintas de la misma colección.

Las trazas de las potencias del transporte determinan además una tabla
periódica que reconstruye el logaritmo del cociente constitutivo.
Su serie de momentos, la representación de Mellin y las cotas de resto
se demuestran en la sección~\ref{amp-afin-sec}. La composición diagonal
y la tensorial mantienen explícita la profundidad que cada operación
conserva; el retorno espectral permanece distinto del retorno del estado.

El teorema~\ref{iii:thm:restitucion-funcional-catalan} completa la
relación mediante una restitución integral: la función constitutiva
fija el momento orientado entero y la condición de límite nulo en el
origen determina su normalización. La prueba distingue la recuperación
de esa colección de la inversión de sus dos evaluaciones de canal.
Ambas operaciones conservan la lectura de velocidad
$\widehat c=(\det\mathcal V)^{-1}$ y su traza normalizada, que enlazan
el transporte constitutivo con la duración y la longitud del mismo estado.

La composición con Barbero--Immirzi conserva una información adicional:
el signo relativo de los canales. La incidencia hexada--octada y el
calendario fijan los términos del funcional angular; la inversa de
$f$ permite expresarlo sobre $\mathcal V$. El
teorema~\ref{iii:thm:barbero-factorization} demuestra esa
factorización. Un intercambio de canales respecto de una marca de
orientación fija invierte el signo del funcional; una conjugación
simultánea de operador y marca lo conserva. De este modo, la respuesta
constitutiva transporta la información que requiere su lectura angular.

% BEGIN SINTESIS_ADITIVA_20260922
\par
La sección~\ref{delta22:iii:restitucion} reúne el potencial espectral de la respuesta constitutiva, la restitución global de los canales mediante velocidad normalizada y Barbero--Immirzi, y sus cotas de estabilidad. Sobre la imagen generada, con las etiquetas y bases conservadas, la composición recupera el par angular, la sección de acción y la lectura periódica del registro dodecafásico.
% END SINTESIS_ADITIVA_20260922

La sección~\ref{iii:sec:incidencias-espectrales} hace explícitas las
normalizaciones de esta composición. La factorización positiva separa
el factor común de la respuesta y su componente unimodular orientada;
ambos datos permiten restituir las cuatro coordenadas constitutivas.
La representación de Fourier distingue después los modos de las
incidencias $90/120$, con los pesos firmados de la cadena y una
normalización de transformada declarada. Esta organización permite
comparar las realizaciones antecedentes con el operador utilizado
en el artículo, conservando en cada una su función y su dominio.

El apéndice~\ref{iii:hist:antecedentes} desarrolla además la realización
finita sobre posiciones, modos y fase: prueba su equivarianza y la
descomposición en veintisiete componentes, conserva las fases de sus
autovalores y evalúa su sensibilidad paramétrica. Los lectores angulares
antecedentes se reproducen con sus operaciones y datos declarados.
Su comparación con los lectores constitutivos del cuerpo distingue los
pesos, las normalizaciones y los dominios de cada realización.

\subsection{Polarización, carga y realización dimensional}

El desarrollo eléctrico se apoya también en dos resultados de las
hojas aritméticas. El desajuste entre capacidades nonádica y decimal
produce una discrepancia acumulada uniformemente acotada; su diferencia
temporal define el incremento de polarización. Por otra parte, el
emparejamiento de las fibras aditiva y multiplicativa en $D_9^3$
determina el modo trítico y la norma $\sqrt3/4$ del conmutador de sus
proyecciones. La sección~\ref{iii:sec:vacancias-prefactor} conserva
las pruebas y los dominios de ambas construcciones. Su relación con
el límite $3/4$ de la respuesta temporal precisa qué dato comparten
sin identificar sus espacios operatorios.

La sección~\ref{iii:sec:carga} compone la impedancia, $\alpha$ y la
constante de Planck mediante la única sección positiva que satisface
$Z_{\rm vac}e_a^2=2\alpha h_a$. De ella se obtienen resistencia de
von Klitzing, cuanto de conductancia, flujo magnético y constante de
Josephson, con identidades cruzadas exactas. El retorno de la acción
deja invariantes las dos primeras y transforma las dos últimas con
potencias opuestas del cociente de acción. Este contenido eléctrico
es una consecuencia de las secciones anteriores, no una lista de
constantes independientes.

La evaluación final mantiene las rectas de acción, carga y velocidad,
así como las transformaciones de sus coordenadas al cambiar unidades.
La tabla de la sección~\ref{iii:sec:evaluacion} presenta las coordenadas
internas obtenidas por evaluación. Una comparación física se formula después, con la coordenatización
dimensional y el régimen de acoplamiento declarados. El argumento
queda organizado por una cuestión común: qué estructura determina la
respuesta, qué ecuaciones obliga a satisfacer y qué información puede
recuperarse de los valores así obtenidos.
